Under a strict inter-patient protocol, with every choice fixed in advance, class-conditional generative augmentation did not detectably improve rare-arrhythmia detection on MIT-BIH (0 of 4 confirmatory comparisons in its favour), and the follow-up analyses did not change that answer. The synthetic beats came from the same few training patients as the real ones, so they added beats but not the variety of new patients. What did move the score uses the same cues a cardiologist uses: each beat seen against the patient's own normal beat (0.088 to 0.371), prematurity judged against the recording's mean heart period, and no supraventricular calls during atrial fibrillation; with ten runs averaged the pipeline reached 0.587, a score that is optimistic after repeated reads of the test records and lower on new patients. The system delivers what the course asked for, a class, a saliency map and a checked note for every beat; it is a course project and not a medical device. Larger models, more runs and above all more training patients with supraventricular and fusion beats are needed to establish how much augmentation can add.
